% GENERATED FILE. Edit canonical lessons, then run npm run generate:books.
% canonical-source-hash: fnv1a64:68d360ad85d04fe6
% canonical-lessons: TA-C431-varru, TA-C431-ular, TA-C431-aviyaku, TA-C431-karuku, TA-C431-nirmattam, TA-R431-review-424-431

\chapter{Drying Up}
\label{ch:ta-drying-up}
\hlchaptermodality{431}

\begin{chapteropening}
\textbf{By the end of this chapter:} \emph{I can talk about water drying up and evaporating, things getting scorched, and the water level.}

Complete the last lesson of chapter 431: I can talk about water drying up and evaporating, things getting scorched, and the water level.
\end{chapteropening}

\section[\texorpdfstring{\ta{வற்று} (vaṟṟu)}{வற்று (vaṟṟu)}]{\ta{வற்று} (vaṟṟu) — to dry up; to ebb}

\label{lesson:TA-C431-varru}



\begin{quote}
Before the new one: say the Tamil for to glitter, then the Tamil for to sound.
\end{quote}

\subsection*{You'll want to know: \ta{வற்று}}
\textbf{\ta{வற்று}} — \emph{vaṟṟu} — \textquotedblleft{}to dry up; to ebb\textquotedblright{}.

\begin{quote}
\ta{கோடையில்} \ta{ஏரி} \ta{வற்றியது}.
\end{quote}

\begin{quote}
\emph{kōṭaiyil ēri vaṟṟiyatu.}
\end{quote}

\begin{quote}
In summer the lake dried up.
\end{quote}

Past (it …): \ta{வற்றியது}, \emph{vaṟṟiyatu}.

Of a river or a well; \ta{வறட்சி}, a drought, is a close cousin.

\begin{cousinweb}
\textbf{\ta{வற்று}} (\emph{vaṟṟu}) is native Dravidian (DEDR 5320), \textquotedblleft{}to grow dry, to evaporate, to ebb\textquotedblright{}.
\end{cousinweb}

\begin{grammarlens}[title={What you've built}]
One more word for getting around.
\end{grammarlens}

\subsection*{Your turn}
\begin{itemize}
\raggedright
  \item \emph{Say it:} \emph{vaṟṟu}
  \item \emph{Say it:} \emph{vaṟṟu}, once more
  \item \emph{Say it:} \emph{kōṭaiyil ēri vaṟṟiyatu.} — In summer the lake dried up.
  \item \emph{Recall:} say the Tamil for to glitter, then the Tamil for to sound, then say \emph{vaṟṟu} again
\end{itemize}

\begin{tcolorbox}[breakable,colback=teal!4,colframe=teal!35!black,title={Before you move on}]
What does \ta{வற்று} mean? (\textquotedblleft{}To dry up; to ebb\textquotedblright{}.) Say it once more.
\end{tcolorbox}
\section[\texorpdfstring{\ta{உலர்} (ular)}{உலர் (ular)}]{\ta{உலர்} (ular) — to become dry; to dry}

\label{lesson:TA-C431-ular}



\begin{quote}
Before the new one: say the Tamil for to sound, then the Tamil for to dry up.
\end{quote}

\subsection*{You'll want to know: \ta{உலர்}}
\textbf{\ta{உலர்}} — \emph{ular} — \textquotedblleft{}to become dry; to dry\textquotedblright{}.

\begin{quote}
\ta{வெயிலில்} \ta{துணிகள்} \ta{சீக்கிரம்} \ta{உலர்ந்தன}.
\end{quote}

\begin{quote}
\emph{veyilil tuṇikaḷ cīkkiram ularntaṉa.}
\end{quote}

\begin{quote}
The clothes dried quickly in the sun.
\end{quote}

Past (it …): \ta{உலர்ந்தது}, \emph{ularntatu}.

To dry something is \ta{உலர்த்து}.

\begin{cousinweb}
\textbf{\ta{உலர்}} (\emph{ular}) is native Dravidian (DEDR 674), \textquotedblleft{}to become dry, to wither\textquotedblright{}.
\end{cousinweb}

\begin{grammarlens}[title={What you've built}]
One more word for getting around.
\end{grammarlens}

\subsection*{Your turn}
\begin{itemize}
\raggedright
  \item \emph{Say it:} \emph{ular}
  \item \emph{Say it:} \emph{ular}, once more
  \item \emph{Say it:} \emph{veyilil tuṇikaḷ cīkkiram ularntaṉa.} — The clothes dried quickly in the sun.
  \item \emph{Recall:} say the Tamil for to sound, then the Tamil for to dry up, then say \emph{ular} again
\end{itemize}

\begin{tcolorbox}[breakable,colback=teal!4,colframe=teal!35!black,title={Before you move on}]
What does \ta{உலர்} mean? (\textquotedblleft{}To become dry; to dry\textquotedblright{}.) Say it once more.
\end{tcolorbox}
\section[\texorpdfstring{\ta{ஆவியாகு} (āviyāku)}{ஆவியாகு (āviyāku)}]{\ta{ஆவியாகு} (āviyāku) — to evaporate}

\label{lesson:TA-C431-aviyaku}



\begin{quote}
Before the new one: say the Tamil for to dry up, then the Tamil for to become dry.
\end{quote}

\subsection*{You'll want to know: \ta{ஆவியாகு}}
\textbf{\ta{ஆவியாகு}} — \emph{āviyāku} — \textquotedblleft{}to evaporate\textquotedblright{}.

\begin{quote}
\ta{வெயிலில்} \ta{தண்ணீர்} \ta{ஆவியானது}.
\end{quote}

\begin{quote}
\emph{veyilil taṇṇīr āviyāṉatu.}
\end{quote}

\begin{quote}
The water evaporated in the sun.
\end{quote}

Past (it …): \ta{ஆவியானது}, \emph{āviyāṉatu}.

\ta{ஆவி}, steam, + \ta{ஆகு}, to become.

\begin{cousinweb}
\textbf{\ta{ஆவியாகு}} (\emph{āviyāku}) is built from \ta{ஆவி}, \textquotedblleft{}breath, steam\textquotedblright{} (DEDR 393), and \ta{ஆகு}, \textquotedblleft{}to become\textquotedblright{}.
\end{cousinweb}

\begin{grammarlens}[title={What you've built}]
One more word for getting around.
\end{grammarlens}

\subsection*{Your turn}
\begin{itemize}
\raggedright
  \item \emph{Say it:} \emph{āviyāku}
  \item \emph{Say it:} \emph{āviyāku}, once more
  \item \emph{Say it:} \emph{veyilil taṇṇīr āviyāṉatu.} — The water evaporated in the sun.
  \item \emph{Recall:} say the Tamil for to dry up, then the Tamil for to become dry, then say \emph{āviyāku} again
\end{itemize}

\begin{tcolorbox}[breakable,colback=teal!4,colframe=teal!35!black,title={Before you move on}]
What does \ta{ஆவியாகு} mean? (\textquotedblleft{}To evaporate\textquotedblright{}.) Say it once more.
\end{tcolorbox}
\section[\texorpdfstring{\ta{கருகு} (karuku)}{கருகு (karuku)}]{\ta{கருகு} (karuku) — to get scorched; to be charred}

\label{lesson:TA-C431-karuku}



\begin{quote}
Before the new one: say the Tamil for to become dry, then the Tamil for to evaporate.
\end{quote}

\subsection*{You'll want to know: \ta{கருகு}}
\textbf{\ta{கருகு}} — \emph{karuku} — \textquotedblleft{}to get scorched; to be charred\textquotedblright{}.

\begin{quote}
\ta{கடும்} \ta{வெயிலில்} \ta{செடிகள்} \ta{கருகின}.
\end{quote}

\begin{quote}
\emph{kaṭum veyilil ceṭikaḷ karukiṉa.}
\end{quote}

\begin{quote}
The plants were scorched in the fierce sun.
\end{quote}

Past (it …): \ta{கருகியது}, \emph{karukiyatu}.

Of food: \ta{சாதம்} \ta{கருகியது}, the rice got burnt.

\begin{cousinweb}
\textbf{\ta{கருகு}} (\emph{karuku}) is native Dravidian (DEDR 1278), \textquotedblleft{}to be scorched, blackened by fire or sun\textquotedblright{}, the family of \ta{கரு}, \textquotedblleft{}black\textquotedblright{}.
\end{cousinweb}

\begin{grammarlens}[title={What you've built}]
One more word for getting around.
\end{grammarlens}

\subsection*{Your turn}
\begin{itemize}
\raggedright
  \item \emph{Say it:} \emph{karuku}
  \item \emph{Say it:} \emph{karuku}, once more
  \item \emph{Say it:} \emph{kaṭum veyilil ceṭikaḷ karukiṉa.} — The plants were scorched in the fierce sun.
  \item \emph{Recall:} say the Tamil for to become dry, then the Tamil for to evaporate, then say \emph{karuku} again
\end{itemize}

\begin{tcolorbox}[breakable,colback=teal!4,colframe=teal!35!black,title={Before you move on}]
What does \ta{கருகு} mean? (\textquotedblleft{}To get scorched; to be charred\textquotedblright{}.) Say it once more.
\end{tcolorbox}
\section[\texorpdfstring{\ta{நீர்மட்டம்} (nīrmaṭṭam)}{நீர்மட்டம் (nīrmaṭṭam)}]{\ta{நீர்மட்டம்} (nīrmaṭṭam) — the water level}

\label{lesson:TA-C431-nirmattam}



\begin{quote}
Before the new one: say the Tamil for to evaporate, then the Tamil for to get scorched.
\end{quote}

\subsection*{You'll want to know: \ta{நீர்மட்டம்}}
\textbf{\ta{நீர்மட்டம்}} — \emph{nīrmaṭṭam} — \textquotedblleft{}the water level\textquotedblright{}.

\begin{quote}
\ta{அணையின்} \ta{நீர்மட்டம்} \ta{குறைந்தது}.
\end{quote}

\begin{quote}
\emph{aṇaiyiṉ nīrmaṭṭam kuṟaintatu.}
\end{quote}

\begin{quote}
The water level in the dam fell.
\end{quote}

\ta{நீர்}, water, + \ta{மட்டம்}, a level.

\begin{cousinweb}
\textbf{\ta{நீர்மட்டம்}} (\emph{nīrmaṭṭam}) is built from \ta{நீர்}, \textquotedblleft{}water\textquotedblright{} (DEDR 3690), and \ta{மட்டம்}, \textquotedblleft{}a level, a measure\textquotedblright{} (DEDR 4660).
\end{cousinweb}

\begin{grammarlens}[title={What you've built}]
One more word for getting around.
\end{grammarlens}

\subsection*{Your turn}
\begin{itemize}
\raggedright
  \item \emph{Say it:} \emph{nīrmaṭṭam}
  \item \emph{Say it:} \emph{nīrmaṭṭam}, once more
  \item \emph{Say it:} \emph{aṇaiyiṉ nīrmaṭṭam kuṟaintatu.} — The water level in the dam fell.
  \item \emph{Recall:} say the Tamil for to evaporate, then the Tamil for to get scorched, then say \emph{nīrmaṭṭam} again
\end{itemize}

\begin{tcolorbox}[breakable,colback=teal!4,colframe=teal!35!black,title={Before you move on}]
What does \ta{நீர்மட்டம்} mean? (\textquotedblleft{}The water level\textquotedblright{}.) Say it once more.
\end{tcolorbox}
\section[(dialogue)]{Looking back: from Statues and Monuments to Drying Up}

\label{lesson:TA-R431-review-424-431}



\begin{quote}
Say the Tamil for a town square and a statue.
\end{quote}

\subsection*{Your turn}
Say these aloud:

\begin{itemize}
\raggedright
  \item a town square, a statue, a monument, a mansion, an embassy
  \item a lighthouse, a planetarium, a site, a view, a canteen
  \item bustle, noise, to gather, to fill up, to gasp
  \item climate, warmth, temperature, to scorch, to be sultry
  \item a spray of rain, to drizzle, to pour down, to stagnate, to drain away
  \item a forecast, a cyclone, hailstones, a downpour, cloud cover
  \item sunrise, brightness, the horizon, to glitter, to sound
  \item to dry up, to become dry, to evaporate, to get scorched, the water level
\end{itemize}

\begin{tcolorbox}[breakable,colback=teal!4,colframe=teal!35!black,title={Before you move on}]
A town square? (\textbf{\ta{சதுக்கம்}}, \emph{catukkam}.) A statue? (\textbf{\ta{சிலை}}, \emph{cilai}.) A monument? (\textbf{\ta{நினைவுச்சின்னம்}}, \emph{niṉaivucciṉṉam}.) A mansion? (\textbf{\ta{மாளிகை}}, \emph{māḷikai}.) An embassy? (\textbf{\ta{தூதரகம்}}, \emph{tūtarakam}.) A lighthouse? (\textbf{\ta{கலங்கரை} \ta{விளக்கம்}}, \emph{kalaṅkarai viḷakkam}.) A planetarium? (\textbf{\ta{கோளரங்கம்}}, \emph{kōḷaraṅkam}.) A site? (\textbf{\ta{தலம்}}, \emph{talam}.) A view? (\textbf{\ta{காட்சி}}, \emph{kāṭci}.) A canteen? (\textbf{\ta{சிற்றுண்டிச்சாலை}}, \emph{ciṟṟuṇṭiccālai}.) Bustle? (\textbf{\ta{சந்தடி}}, \emph{cantaṭi}.) Noise? (\textbf{\ta{இரைச்சல்}}, \emph{iraiccal}.) To gather? (\textbf{\ta{திரள்}}, \emph{tiraḷ}.) To fill up? (\textbf{\ta{நிரம்பு}}, \emph{nirampu}.) To gasp? (\textbf{\ta{திணறு}}, \emph{tiṇaṟu}.) Climate? (\textbf{\ta{காலநிலை}}, \emph{kālanilai}.) Warmth? (\textbf{\ta{வெப்பம்}}, \emph{veppam}.) Temperature? (\textbf{\ta{வெப்பநிலை}}, \emph{veppanilai}.) To scorch? (\textbf{\ta{கொளுத்து}}, \emph{koḷuttu}.) To be sultry? (\textbf{\ta{புழுங்கு}}, \emph{puḻuṅku}.) A spray of rain? (\textbf{\ta{சாரல்}}, \emph{cāral}.) To drizzle? (\textbf{\ta{தூறு}}, \emph{tūṟu}.) To pour down? (\textbf{\ta{பொழி}}, \emph{poḻi}.) To stagnate? (\textbf{\ta{தேங்கு}}, \emph{tēṅku}.) To drain away? (\textbf{\ta{வடி}}, \emph{vaṭi}.) A forecast? (\textbf{\ta{முன்னறிவிப்பு}}, \emph{muṉṉaṟivippu}.) A cyclone? (\textbf{\ta{சூறாவளி}}, \emph{cūṟāvaḷi}.) Hailstones? (\textbf{\ta{ஆலங்கட்டி}}, \emph{ālaṅkaṭṭi}.) A downpour? (\textbf{\ta{அடைமழை}}, \emph{aṭaimaḻai}.) Cloud cover? (\textbf{\ta{மேகமூட்டம்}}, \emph{mēkamūṭṭam}.) Sunrise? (\textbf{\ta{உதயம்}}, \emph{utayam}.) Brightness? (\textbf{\ta{வெளிச்சம்}}, \emph{veḷiccam}.) The horizon? (\textbf{\ta{அடிவானம்}}, \emph{aṭivāṉam}.) To glitter? (\textbf{\ta{மின்னு}}, \emph{miṉṉu}.) To sound? (\textbf{\ta{ஒலி}}, \emph{oli}.) To dry up? (\textbf{\ta{வற்று}}, \emph{vaṟṟu}.) To become dry? (\textbf{\ta{உலர்}}, \emph{ular}.) To evaporate? (\textbf{\ta{ஆவியாகு}}, \emph{āviyāku}.) To get scorched? (\textbf{\ta{கருகு}}, \emph{karuku}.) The water level? (\textbf{\ta{நீர்மட்டம்}}, \emph{nīrmaṭṭam}.)
\end{tcolorbox}
